\NSPage{079}%
Let \NSi{NS-F-P079-001}{V_m(v,w)} be the forward fundamental matrix for the
homogeneous equation \NSi{NS-F-P079-002}{z'=A_mz}, normalized by
\NSi{NS-F-P079-003}{V_m(w,w)=I_2}; thus
\NSi{NS-F-P079-004}{\partial_vV_m(v,w)=A_m(v)V_m(v,w)}. The scalar nature of
the damping gives the exact identity
\NSFormula{NS-F-P079-005}{display}{7.18}
\begin{equation}
  V_m(v,w)=\exp\!\left(-(m^2-1)\int_w^v d(a)\,da\right)V_1(v,w).
  \tag{7.18}\label{eq:7.18}
\end{equation}
The Euclidean norm inequality for the \NSi{NS-F-P079-006}{m=1} equation is
\NSi{NS-F-P079-007}{\partial_v|z|\leq(\lambda-d_{\mathrm{ref}}+C/S_*)|z|}.
As \NSi{NS-F-P079-008}{L_s/S_*} is bounded, it follows that
\NSFormula{NS-F-P079-009}{display}{7.19}
\begin{equation}
  \lVert V_m(v,w)\rVert\leq C P(v)/P(w),
  \qquad 0\leq w\leq v\leq L_s,\qquad m\neq0.
  \tag{7.19}\label{eq:7.19}
\end{equation}
This value bound is uniform in all nonzero integers \NSi{NS-F-P079-010}{m},
since the exponential in \eqref{eq:7.18} is at most one.

\textbf{Step 2: values and coefficient derivatives.}
The zero initial condition gives Duhamel's formula. The bound on
\NSi{NS-F-P079-011}{g_m(w)} contains \NSi{NS-F-P079-012}{P(w)}, which cancels
the denominator in \eqref{eq:7.19}; integration contributes
\NSi{NS-F-P079-013}{L_s=O(S_*)}:
\NSFormula{NS-F-P079-014}{display}{none}
\[
  z_m(v)=\int_0^v V_m(v,w)g_m(w)\,dw,
  \qquad
  |z_m(v)|\leq C\varepsilon^\alpha S_*^{b_0+1}\sqrt{\zeta}\,\delta^{-a_0}P(v),
\]
after enlarging \NSi{NS-F-P079-015}{b_0} for the frame factors. To
differentiate this formula, first use the slow variables, the transverse
coordinate at \NSi{NS-F-P079-016}{v=0}, and the pulse coordinate. In the next
equation, \NSi{NS-F-P079-017}{D^I} differentiates only the slow variables and
initial transverse coordinate, holding \NSi{NS-F-P079-018}{v} fixed. These
parameter derivatives commute with \NSi{NS-F-P079-019}{\partial_v}, so
differentiating the two-dimensional equation gives
\NSFormula{NS-F-P079-020}{display}{7.20}
\begin{equation}
  (D^I z_m)'=A_mD^Iz_m+D^Ig_m+
  \sum_{0<J\leq I}\binom{I}{J}(D^JA_m)D^{I-J}z_m.
  \tag{7.20}\label{eq:7.20}
\end{equation}
All these parameter derivatives of the zero initial data vanish. If the
coefficient derivatives obey
\NSi{NS-F-P079-021}{|D^J A_m|\leq C_JS_*^{c_J}}, an induction using
\eqref{eq:7.19} increases the polynomial degree by at most
\NSFormula{NS-F-P079-022}{display}{none}
\[
  1+\max\{b_I,\max_{0<J\leq I}(c_J+b'_{I-J})\}
\]
at the \NSi{NS-F-P079-023}{I}th step. The inverse-distance exponent can be
enlarged to the maximum of the source and lower-order exponents. The slow
point is fixed in the integral, so \NSi{NS-F-P079-024}{\sqrt{\zeta}} and
\NSi{NS-F-P079-025}{\delta} are fixed there. Pulse-coordinate derivatives
then follow from \eqref{eq:7.17}. This proves every fixed-order derivative
estimate with no decrease of \NSi{NS-F-P079-026}{\alpha}. Differentiation in
the moving frame retains the constraint
\NSi{NS-F-P079-027}{n_\Phi\cdot t_m=0} and accounts for derivatives of the
plane \NSi{NS-F-P079-028}{n_\Phi^\perp}.

\textbf{Step 3: common-torus coordinates and support.}
The preceding estimates used \NSi{NS-F-P079-029}{\xi_g,v} on one lifted
rectangle. To express the derivative estimates in \NSi{NS-F-P079-030}{H}, we
use the directions \NSi{NS-F-P079-031}{v_r,v_t} from \eqref{eq:6.2} and
recall the linear functional
\NSi{NS-F-P079-032}{\lambda_t(w)=v_t\cdot w/(1+b_g^2)}, so
\NSi{NS-F-P079-033}{\lambda_t(v_r)=0} and
\NSi{NS-F-P079-034}{\lambda_t(v_t)=1}. For a current point
\NSi{NS-F-P079-035}{H} with pulse coordinate
\NSi{NS-F-P079-036}{v(H)}, the point on the same path at pulse coordinate
\NSi{NS-F-P079-037}{w} is
\NSFormula{NS-F-P079-038}{display}{none}
\[
  H_w=H+T_g^{-(i-i_0)}c_i\bigl(w-v(H)\bigr)v_t.
\]
Since \NSi{NS-F-P079-039}{v_t} is an eigenvector of
\NSi{NS-F-P079-040}{J_g} with eigenvalue \NSi{NS-F-P079-041}{T_g},
\NSFormula{NS-F-P079-042}{display}{none}
\[
  D_HH_w=I_2-v_t\lambda_t,\qquad
  D_H^aH_w=0\quad(a\geq2),\qquad
  D_Hv=T_g^{i-i_0}\lambda_t/c_i=O(S_*).
\]
Thus the pullback \NSi{NS-F-P079-043}{H\mapsto H_w} at fixed
\NSi{NS-F-P079-044}{w} has bounded derivatives, and returning to the current
common coordinate introduces only powers of \NSi{NS-F-P079-045}{S_*}
multiplying fixed derivatives in the pulse coordinate. Together with Step 2,
this proves \eqref{eq:7.14} in the stated coefficient variables. Deck
translations reindex the copies and commute with this path. Uniqueness of the
initial-value
\NSPage{080}%
equation with zero data proves descent to \NSi{NS-F-P080-001}{H}, without
identifying sources on different preimages of the finer band torus. The path
holds the slow and transverse points fixed. If the source vanishes along their
whole rectangle, the solution does also. The differentiated equation, applied
at their support boundaries, proves smooth zero extension and support
containment.

\textbf{Step 4: the constraint and pressure.}
The definition of \NSi{NS-F-P080-002}{\mathcal A_\Phi} gives
\NSFormula{NS-F-P080-003}{display}{none}
\[
  n_\Phi\cdot t'_m=-n'_\Phi\cdot t_m-m^2d(n_\Phi\cdot t_m),
  \qquad
  (n_\Phi\cdot t_m)'=-m^2d(n_\Phi\cdot t_m).
\]
Zero initial data therefore preserve
\NSi{NS-F-P080-004}{n_\Phi\cdot t_m=0}. Substitution in the projected equation
yields
\NSFormula{NS-F-P080-005}{display}{none}
\[
  t'_m+\mathcal Kt_m+m^2dt_m+f_m
  =\frac{n_\Phi}{|n_\Phi|^2}
  (n_\Phi\cdot\mathcal Kt_m-n'_\Phi\cdot t_m+n_\Phi\cdot f_m).
\]
Multiplication of \eqref{eq:7.13} by
\NSi{NS-F-P080-006}{ikmn_\Phi} gives the negative of this normal vector. Thus
\NSi{NS-F-P080-007}{(t_m,\pi_m)} satisfies the full principal cancellation
\eqref{eq:7.5}. The factor
\NSi{NS-F-P080-008}{(km)^{-1}=O(\varepsilon^{1/2})} and the fixed
coefficient-derivative bounds prove \eqref{eq:7.15}. In particular, normal
components of the source have been canceled as well. On a compact set with
\NSi{NS-F-P080-009}{q>0}, the base has the one-sided endpoint derivative data
of Proposition~\ref{prop:5.5}. The same parameter-dependent equations and
differentiated Duhamel formulas extend those endpoint derivatives from the
source to \NSi{NS-F-P080-010}{t_m} and then to
\NSi{NS-F-P080-011}{\pi_m}.
\end{proof}

The harmonic factor in \eqref{eq:7.18} only adds damping. Consequently the
later appearance of higher harmonics does not force a smaller space-time
domain for these inverses.

\begin{corollary}\label{cor:7.3}
The solution operators \NSi{NS-F-P080-012}{f_m\mapsto(t_m,\pi_m)} for the
projected amplitude equation with \NSi{NS-F-P080-013}{t_m(0)=0} are defined at
every finite correction stage on the same domain
\NSi{NS-F-P080-014}{0<q<q_*}. Increasing the finite harmonic set or the
derivative order changes constants and polynomial degrees, but does not
require shrinking this domain.
\end{corollary}

\begin{proof}
The frame and damping comparisons require only the single finite choice in
Lemma~\ref{lem:7.1}. The exact factor \eqref{eq:7.18} makes every higher
harmonic at least as damped in value as the first. In \eqref{eq:7.20}, the
extra factors of \NSi{NS-F-P080-015}{m} and coefficient derivatives enlarge
only the constants and polynomial degrees in the source bounds. Thus the
previously chosen threshold \NSi{NS-F-P080-016}{q_*} remains valid.
\end{proof}

The zero-data inverse is now available for prescribed harmonic sources. For
the leading stress, we instead need a nonzero real solution whose radial
component stays comparable to \NSi{NS-F-P080-017}{P(v)}. We choose its initial
value in the growing coordinate of the frame \NSi{NS-F-P080-018}{B}; its
amplitude and direction must remain controlled through the subsequent decay.

\begin{lemma}\label{lem:7.4}
For each sign \NSi{NS-F-P080-019}{\sigma}, let
\NSi{NS-F-P080-020}{t^h_\sigma=B(z_+,z_-)^T} solve
\NSi{NS-F-P080-021}{(t^h_\sigma)'=\mathcal A_\Phi t^h_\sigma-dt^h_\sigma} with
\NSi{NS-F-P080-022}{z_+(0)=P(0)}, \NSi{NS-F-P080-023}{z_-(0)=0}. This
specifies the real amplitude of the homogeneous \NSi{NS-F-P080-024}{m=1}
problem, with pressure given by \eqref{eq:7.13} at
\NSi{NS-F-P080-025}{f_m=0}. With
\NSi{NS-F-P080-026}{K_a,N_a,s_a} as in \eqref{eq:7.7}, write
\NSi{NS-F-P080-027}{t^h_\sigma=x(e_r-s_aK_a)+yN_a}. Then
\NSFormula{NS-F-P080-028}{display}{7.21}
\begin{equation}
  cP(v)\leq x(v)\leq CP(v),
  \qquad
  \frac{y(v)}{x(v)}=c_0\sqrt{1+s(v)^2}+O(S_*^{-1}).
  \tag{7.21}\label{eq:7.21}
\end{equation}
Every fixed derivative of this homogeneous solution in the coefficient
variables or pulse coordinate is bounded by
\NSi{NS-F-P080-029}{C_IS_*^{b_I}P(v)}. Its homogeneous pressure, given by
\eqref{eq:7.13} with \NSi{NS-F-P080-030}{f=0}, has the extra factor
\NSi{NS-F-P080-031}{\varepsilon^{1/2}}.
\end{lemma}
\NSPage{081}%
\begin{proof}
We control the ratio of the decaying and growing frame coordinates, then
convert to the radial and tangential coordinates \NSi{NS-F-P081-001}{x,y}.
Write \NSi{NS-F-P081-002}{E=(E_{ab})}. While
\NSi{NS-F-P081-003}{z_+>0}, the ratio
\NSi{NS-F-P081-004}{r=z_-/z_+} satisfies
\NSFormula{NS-F-P081-005}{display}{none}
\[
  r'=E_{21}+(-2\lambda+E_{22}-E_{11})r-E_{12}r^2,
  \qquad r(0)=0.
\]
Let \NSi{NS-F-P081-006}{c_\lambda>0} be a uniform lower bound for
\NSi{NS-F-P081-007}{\lambda}. For a sufficiently large fixed
\NSi{NS-F-P081-008}{K_b} and all sufficiently large
\NSi{NS-F-P081-009}{S_*},
\NSFormula{NS-F-P081-010}{display}{none}
\[
  \left.\pm r'\right|_{r=\pm K_b/S_*}
  \leq
  \frac{-2c_\lambda K_b+C(1+K_b/S_*+K_b^2/S_*^2)}{S_*}<0.
\]
Thus \NSi{NS-F-P081-011}{|r|\leq K_b/S_*}. Since
\NSi{NS-F-P081-012}{z_+(0)=P(0)>0}, the envelope definition gives
\NSFormula{NS-F-P081-013}{display}{none}
\[
  \left(\log\frac{z_+}{P}\right)'
  =-(d-d_{\mathrm{ref}})+E_{11}+E_{12}r=O(S_*^{-1}),
  \qquad
  \left|\log\frac{z_+(v)}{P(v)}\right|
  \leq\frac{Cv}{S_*}\leq\frac{CL_s}{S_*}\leq C.
\]
This prevents \NSi{NS-F-P081-014}{z_+} from reaching zero and proves
\NSi{NS-F-P081-015}{cP\leq z_+\leq CP} on the whole interval.
Multiplication by the matrix in \eqref{eq:7.8} gives
\NSi{NS-F-P081-016}{x=z_++z_-=z_+(1+r)} and
\NSi{NS-F-P081-017}{y=c_0\sqrt{1+s^2}(z_+-z_-)
=c_0\sqrt{1+s^2}\,z_+(1-r)}. These identities give \eqref{eq:7.21}.
Differentiating the coordinate equation and using \eqref{eq:7.19} proves the
derivative estimates just as in \eqref{eq:7.20}; the initial value
\NSi{NS-F-P081-018}{P(0)} is fixed with the label. The pressure formula
supplies the extra factor \NSi{NS-F-P081-019}{\varepsilon^{1/2}}.
\end{proof}

\textbf{Energy transfer from the shear.}
The pulse equation also identifies the source of the homogeneous growth. For
the real amplitude \NSi{NS-F-P081-020}{t=t^h_\sigma}, the constraint
\NSi{NS-F-P081-021}{n_\Phi\cdot t=0} makes the normal term in
\NSi{NS-F-P081-022}{\mathcal A_\Phi t} orthogonal to \NSi{NS-F-P081-023}{t}.
Moreover,
\NSFormula{NS-F-P081-024}{display}{none}
\[
  t\cdot\mathcal Kt=-2F t_rt_\theta+(2F+RF_R)t_rt_\theta+G_Rt_rt_z
  =g\cdot\bigl(t_r(t_\theta,t_z)\bigr).
\]
Taking the scalar product of
\NSi{NS-F-P081-025}{t'=\mathcal A_\Phi t-dt} with \NSi{NS-F-P081-026}{t} gives
\NSFormula{NS-F-P081-027}{display}{7.22}
\begin{equation}
  \frac12\frac{d}{dv}|t|^2
  =-g\cdot\bigl(t_r(t_\theta,t_z)\bigr)
  -\varepsilon k^2|n_\Phi|^2|t|^2.
  \tag{7.22}\label{eq:7.22}
\end{equation}
The first term gives energy exchange with the base shear; the second gives
viscous dissipation. The vector
\NSi{NS-F-P081-028}{t_r(t_\theta,t_z)} records transport of azimuthal and
axial momentum across a cylinder. At the representative point, decompose such
a flux vector as \NSi{NS-F-P081-029}{T=T_NN+T_KK}. The energy transfer is
\NSi{NS-F-P081-030}{-g_0\cdot T=-|g_0|T_N}; the covariance construction also
prescribes the transverse component \NSi{NS-F-P081-031}{T_K}. We therefore
need both components of the covariance to match the prescribed momentum flux.

Viscosity remains a leading-order term in the pulse equation. At viscosity
one, the base velocity and radial length give Reynolds number of order
\NSi{NS-F-P081-032}{q^{-h}}. The carrier wavelength is
\NSi{NS-F-P081-033}{\sqrt{Q}/k\asymp Q^{1/2+h/2}}, whereas the physical
amplitude of the leading oscillations constructed below has order
\NSi{NS-F-P081-034}{Q^{-A}\sqrt{\varepsilon}=Q^{-1/2-h/2}}, up to powers of
\NSi{NS-F-P081-035}{S_*}. The powers of \NSi{NS-F-P081-036}{Q} in their
product cancel. Thus \NSi{NS-F-P081-037}{\varepsilon k^2\asymp1} in the pulse
equation throughout the correction process.

\subsection{Covariance of the leading oscillations and its linearization}
\label{sec:7.3}
The two real homogeneous pulses \NSi{NS-F-P081-038}{t^h_\pm} from
Lemma~\ref{lem:7.4} determine two covariance directions in each slow box.
Their scalar amplitudes \NSi{NS-F-P081-039}{a_\pm} will be chosen in
Proposition~\ref{prop:7.5} so that a positive combination of these directions
equals the leading stress. In this subsection we will compute the two
covariances, solve for the squared amplitudes, and then differentiate this
relation to obtain signed stress corrections using
\NSi{NS-F-P081-040}{\mathcal L} from Proposition~\ref{prop:7.6}.

Recall the angular and auxiliary averages in \eqref{eq:3.7}, with normalized
Haar measure on the auxiliary torus. They are taken at fixed slow variables,
before evaluation at the physical phase map. For real velocity fields
\NSi{NS-F-P081-041}{w,v} on the extended domain, we write
\NSi{NS-F-P081-042}{w_{\mathrm{tan}}=(w_\theta,w_z)}.
\NSPage{082}%
The product \NSi{NS-F-P082-001}{w_rw_{\mathrm{tan}}} records radial transport
of azimuthal and axial momentum. At each fixed slow point, we define the
covariance \NSi{NS-F-P082-002}{\mathcal C(w)} and symmetric cross term
\NSi{NS-F-P082-003}{\mathcal B(w,v)} by
\NSFormula{NS-F-P082-004}{display}{7.23}
\begin{equation}
  \mathcal C(w)=\bigl\langle\langle w_rw_{\mathrm{tan}}\rangle_\theta\bigr\rangle_Y,
  \qquad
  \mathcal B(w,v)=\bigl\langle\langle
  w_rv_{\mathrm{tan}}+v_rw_{\mathrm{tan}}
  \rangle_\theta\bigr\rangle_Y.
  \tag{7.23}\label{eq:7.23}
\end{equation}
The averages can be computed on any applicable common torus by
Lemma~\ref{lem:6.2}. Expansion of the product gives
\NSi{NS-F-P082-005}{D\mathcal C(w)[v]=\mathcal B(w,v)}.

Write \NSi{NS-F-P082-006}{\beta=(\ell,a)} for a slow box and
\NSi{NS-F-P082-007}{\gamma=(\beta,\sigma)} for one of its two rectangle
labels. The two signs have the same slow cutoff
\NSi{NS-F-P082-008}{\eta_\beta=\eta_{(\beta,+)}=\eta_{(\beta,-)}}; thus the
squared partition sums over \NSi{NS-F-P082-009}{\beta}, with each box counted
once. On the slow neighborhood of \NSi{NS-F-P082-010}{\beta}, we define
\NSFormula{NS-F-P082-011}{display}{none}
\[
  b_\sigma=\chi_g(\xi_g)\psi(v)t^h_\sigma\cos(k\Phi_\sigma),
  \qquad \sigma\in\{+,-\},
\]
and extend by zero off that sign's auxiliary rectangle. The factors
\NSi{NS-F-P082-012}{\chi_g,\psi} give smooth extension there. The field
\NSi{NS-F-P082-013}{b_\sigma} includes neither the slow cutoff
\NSi{NS-F-P082-014}{\eta_\beta} nor a scalar amplitude. The two signs have
disjoint auxiliary supports.

Let \NSi{NS-F-P082-015}{\mathsf H=[\mathcal C(b_+)\mid\mathcal C(b_-)]}, a real
\NSi{NS-F-P082-016}{2\times2} matrix depending on the slow point. Its columns
have component order \NSi{NS-F-P082-017}{(r\theta,rz)}. For any real numbers
\NSi{NS-F-P082-018}{a_+,a_-}, disjointness gives
\NSi{NS-F-P082-019}{\mathcal C(a_+b_++a_-b_-)=\mathsf H(a_+^2,a_-^2)^T}. Thus the
nonnegative span of its columns is exactly the set of stresses obtainable by
these two waves. The strict cone condition will give positive coefficients for
the leading stress, whose square roots define the real wave amplitudes.

\begin{proposition}\label{prop:7.5}
For each slow box \NSi{NS-F-P082-020}{\beta}, use the homogeneous pulses of
Lemma~\ref{lem:7.4} to form \NSi{NS-F-P082-021}{b_\pm} and
\NSi{NS-F-P082-022}{\mathsf H} as above, and use the order-zero target
\NSi{NS-F-P082-023}{\mathcal T_{0,*}=(Q/q)^{A+1/2}\mathcal T_0}. After a further single decrease
of \NSi{NS-F-P082-024}{q_*}, the matrix \NSi{NS-F-P082-025}{\mathsf H} is invertible
throughout each enlarged slow neighborhood, and the components of
\NSi{NS-F-P082-026}{\mathsf H^{-1}\mathcal T_{0,*}} are positive on its intersection with
\NSi{NS-F-P082-027}{X_a<X<X_b}. We may therefore define, with componentwise
square roots,
\NSFormula{NS-F-P082-028}{display}{7.24}
\begin{equation}
  \boldsymbol y=\mathsf H^{-1}\mathcal T_{0,*},
  \qquad a_\sigma=\sqrt{y_\sigma},
  \qquad W_0=\sqrt{\varepsilon}\sum_{\sigma=\pm}a_\sigma b_\sigma.
  \tag{7.24}\label{eq:7.24}
\end{equation}
The squared amplitudes satisfy, for every fixed slow multi-index
\NSi{NS-F-P082-029}{I} on \NSi{NS-F-P082-030}{X_a<X<X_b},
\NSFormula{NS-F-P082-031}{display}{7.25}
\begin{equation}
  y_\sigma\geq c\sqrt{S_*}\,\zeta,
  \qquad
  |D^Iy_\sigma|\leq C_IS_*^{b_I}\zeta\delta^{-a_I}.
  \tag{7.25}\label{eq:7.25}
\end{equation}
The coefficients \NSi{NS-F-P082-032}{a_\sigma} extend smoothly by zero at the
shell edges. On the label's slow neighborhood, the harmonic coefficients of
\NSi{NS-F-P082-033}{W_0} obey the derivative and envelope estimates of
\NSi{NS-F-P082-034}{\mathsf W_{1/2}}. Multiplication by the slow cutoff gives
\NSi{NS-F-P082-035}{\eta_\beta W_0\in\mathsf W_{1/2}}. Before that multiplication,
the covariance is
\NSFormula{NS-F-P082-036}{display}{7.26}
\begin{equation}
  \mathcal C(W_0)=\varepsilon\mathcal T_{0,*}.
  \tag{7.26}\label{eq:7.26}
\end{equation}
This positive-coefficient decomposition uses only the order-zero target
\NSi{NS-F-P082-037}{\mathcal T_{0,*}}.
\end{proposition}

\begin{proof}
The proof of Lemma~\ref{lem:7.1} applies throughout these neighborhoods: their
diameter remains \NSi{NS-F-P082-038}{O(S_*^{-3})}, and the base comparison
holds on the enlarged annulus. The phase and homogeneous-pulse estimates, and
hence the column estimates below, hold there with uniform constants.

\textbf{Step 1: compute and estimate the covariance columns.}
The angular frequency \NSi{NS-F-P082-039}{kp} is a nonzero integer, so the
angular average of \NSi{NS-F-P082-040}{\cos^2(k\Phi_\sigma)} is exactly
\NSi{NS-F-P082-041}{1/2}. We may compute the auxiliary average on the band
torus. There the rectangle parametrization has area element
\NSi{NS-F-P082-042}{|\det(v_r,v_t)|\,d\xi_g\,d\eta_g}, with
\NSi{NS-F-P082-043}{d\eta_g=c_i\,dv}. Recall that
\NSi{NS-F-P082-044}{x,y} are the geometric coordinates of
\NSi{NS-F-P082-045}{t^h_\sigma} in Lemma~\ref{lem:7.4}. This amplitude is
independent of \NSi{NS-F-P082-046}{\xi_g} and has radial component
\NSi{NS-F-P082-047}{x}. Therefore the column
\NSi{NS-F-P082-048}{\mathsf H_\sigma} equals
\NSFormula{NS-F-P082-049}{display}{7.27}
\begin{equation}
  \mathsf H_\sigma=
  \frac{|\det(v_r,v_t)|}{2}
  \left(\int\chi_g^2\,d\xi_g\right)c_i
  \int_0^{L_s}\psi^2x\,t^h_{\sigma,\mathrm{tan}}\,dv.
  \tag{7.27}\label{eq:7.27}
\end{equation}
\NSPage{083}%
Define the corresponding positive scalar by
\NSFormula{NS-F-P083-001}{display}{none}
\[
  h_\sigma=
  \frac{|\det(v_r,v_t)|}{2}
  \left(\int\chi_g^2\,d\xi_g\right)c_i
  \int_0^{L_s}\psi^2x^2\,dv.
\]
The Gaussian envelope and the region where \NSi{NS-F-P083-002}{\psi=1} give
\NSFormula{NS-F-P083-003}{display}{none}
\[
  c\sqrt{L_s}\leq\int_0^{L_s}\psi^2x^2\,dv\leq C\sqrt{L_s},
  \qquad
  \int_0^{L_s}\frac{|v-L_s/2|}{L_s}\psi^2x^2\,dv\leq C.
\]
Indeed the substitution
\NSi{NS-F-P083-004}{w=(v-L_s/2)/\sqrt{L_s}} reduces these to integrals of
\NSi{NS-F-P083-005}{e^{-cw^2}} and
\NSi{NS-F-P083-006}{|w|e^{-cw^2}}. Hence
\NSi{NS-F-P083-007}{h_\sigma\asymp c_i\sqrt{L_s}\asymp S_*^{-1/2}}, and the
normalized \NSi{NS-F-P083-008}{x^2\psi^2} average of
\NSi{NS-F-P083-009}{|v-L_s/2|/L_s} is
\NSi{NS-F-P083-010}{O(S_*^{-1/2})}. The definitions of
\NSi{NS-F-P083-011}{K_a,N_a,s_a} in \eqref{eq:7.7} give the tangential part
of the pulse:
\NSi{NS-F-P083-012}{t^h_{\sigma,\mathrm{tan}}=-s_axK_a+yN_a}. By
Lemma~\ref{lem:7.4}, its ratio to \NSi{NS-F-P083-013}{x} differs by
\NSi{NS-F-P083-014}{O(S_*^{-1})} from
\NSi{NS-F-P083-015}{-s(v)K+c_0\sqrt{1+s(v)^2}N}. The preceding weighted
average concentrates near \NSi{NS-F-P083-016}{v=L_s/2}, where
\NSi{NS-F-P083-017}{s=\sigma u_*}. Thus, setting
\NSi{NS-F-P083-018}{A_c=-c_0\sqrt{1+u_*^2}>0}, we obtain
\NSFormula{NS-F-P083-019}{display}{7.28}
\begin{equation}
  \mathsf H_\sigma=h_\sigma(-A_cN-\sigma u_*K+e_\sigma),
  \qquad |e_\sigma|\leq CS_*^{-1/2}.
  \tag{7.28}\label{eq:7.28}
\end{equation}
The derivatives of the covariance columns and homogeneous solutions have
polynomial bounds by Lemma~\ref{lem:7.4}.

\textbf{Step 2: solve for positive squared amplitudes.}
First omit the errors \NSi{NS-F-P083-020}{e_\sigma} in \eqref{eq:7.28}. For a
target \NSi{NS-F-P083-021}{T=T_NN+T_KK}, the two scalar equations then give
\NSFormula{NS-F-P083-022}{display}{none}
\[
  h_+y_+=\tfrac12(-T_N/A_c-T_K/u_*),
  \qquad
  h_-y_-=\tfrac12(-T_N/A_c+T_K/u_*).
\]
Thus the reference cone is \NSi{NS-F-P083-023}{T_N<0},
\NSi{NS-F-P083-024}{|T_K|<(u_*/A_c)(-T_N)}. The two formulas above give the
positive coefficient decomposition. Our choice of
\NSi{NS-F-P083-025}{u_*} and the uniform leading cone margin give
\NSi{NS-F-P083-026}{|T_K|/u_*\leq(1-\eta_c)(-T_N/A_c)} for some fixed
\NSi{NS-F-P083-027}{\eta_c>0} and \NSi{NS-F-P083-028}{T=\mathcal T_{0,*}}, by
\eqref{eq:7.1}. Freezing \NSi{NS-F-P083-029}{N,K,c_0} at the representative
changes this normalized cone inequality by an error tending uniformly to
zero. The smooth limiting stress direction gives the same comparison at the
radial edges. The errors in \eqref{eq:7.28} preserve this strict inequality
and give
\NSFormula{NS-F-P083-030}{display}{7.29}
\begin{equation}
  |\det\mathsf H|\geq c/S_*,
  \qquad \lVert\mathsf H^{-1}\rVert\leq C\sqrt{S_*},
  \qquad c\sqrt{S_*}|T|\leq y_\sigma\leq C\sqrt{S_*}|T|.
  \tag{7.29}\label{eq:7.29}
\end{equation}

\textbf{Step 3: derivatives, shell edges, and covariance.}
Apply the leading stress bounds \eqref{eq:4.27} in the chart. Differentiating
\NSi{NS-F-P083-031}{\mathsf H\mathsf H^{-1}=I_2} gives polynomial bounds for all fixed
derivatives of \NSi{NS-F-P083-032}{\mathsf H^{-1}}; differentiating
\NSi{NS-F-P083-033}{\boldsymbol y=\mathsf H^{-1}\mathcal T_{0,*}} then proves \eqref{eq:7.25}. For
\NSi{NS-F-P083-034}{|I|\geq1}, the chain rule gives
\NSFormula{NS-F-P083-035}{display}{none}
\[
  D^Ia_\sigma=
  \sum_{n=1}^{|I|}
  \sum_{\substack{I_1+\cdots+I_n=I\\|I_\nu|\geq1}}
  c_{I_1,\ldots,I_n}y_\sigma^{1/2-n}
  \prod_{\nu=1}^nD^{I_\nu}y_\sigma.
\]
The bounds in \eqref{eq:7.25} therefore imply
\NSFormula{NS-F-P083-036}{display}{none}
\[
  |D^Ia_\sigma|\leq C_IS_*^{b'_I}\sqrt{\zeta}\,\delta^{-a'_I}.
\]
The same bound holds for \NSi{NS-F-P083-037}{I=0}. Its right-hand side tends
to zero at each shell edge for every fixed derivative order, so the zero
extension of \NSi{NS-F-P083-038}{a_\sigma} is smooth. The bounds for
\NSi{NS-F-P083-039}{a_\sigma}, Lemma~\ref{lem:7.4}, and the cutoffs
\NSi{NS-F-P083-040}{\chi_g,\psi} give the local derivative and envelope bounds
for \NSi{NS-F-P083-041}{W_0}. Multiplication by
\NSi{NS-F-P083-042}{\eta_\beta} then gives
\NSi{NS-F-P083-043}{\eta_\beta W_0\in\mathsf W_{1/2}}. Disjointness of the two
rectangles and \eqref{eq:7.24} yield
\NSi{NS-F-P083-044}{\mathcal C(W_0)=\varepsilon\mathsf H\boldsymbol y}, which is \eqref{eq:7.26}.
\end{proof}
\NSPage{084}%
Each slow box now has the required local covariance. Multiplication by the
common slow cutoff \NSi{NS-F-P084-001}{\eta_\beta} multiplies
\eqref{eq:7.26} by \NSi{NS-F-P084-002}{\eta_\beta^2}. Returning to physical
velocities and summing slow boxes and bands gives the leading physical stress
pair
\NSFormula{NS-F-P084-003}{display}{7.30}
\begin{equation}
  \sum_{\beta=(\ell,a)}Q^{-2A}\eta_\beta^2\varepsilon\mathcal T_{0,*}
  =q^{-A-1/2}\mathcal T_0.
  \tag{7.30}\label{eq:7.30}
\end{equation}
To check the scale explicitly, \NSi{NS-F-P084-004}{A=1/2+h} gives
\NSFormula{NS-F-P084-005}{display}{none}
\[
  Q^{-2A}\varepsilon\mathcal T_{0,*}
  =Q^{-2A+h}(Q/q)^{A+1/2}\mathcal T_0
  =Q^{-A+h+1/2}q^{-A-1/2}\mathcal T_0
  =q^{-A-1/2}\mathcal T_0.
\]
We then use \NSi{NS-F-P084-006}{\sum_\beta\eta_\beta^2=1}; cross-label
products vanish by the disjoint auxiliary supports.

The sum supplies the leading physical stress. Later residuals require stress
increments of either sign, which we obtain by varying the amplitudes to first
order about this fixed positive representation. We keep the positive
coefficients \NSi{NS-F-P084-007}{a_\sigma} fixed at their values in
Proposition~\ref{prop:7.5}. Let \NSi{NS-F-P084-008}{\alpha\in\mathbb R}, and
let \NSi{NS-F-P084-009}{\Sigma(R,Z,T)\in\mathsf M_\alpha} be a real two-vector
independent of both angular and auxiliary variables. This auxiliary
independence is an additional requirement here; the class
\NSi{NS-F-P084-010}{\mathsf M_\alpha} alone allows auxiliary dependence. On the
label's slow neighborhood, we define
\NSFormula{NS-F-P084-011}{display}{7.31}
\begin{equation}
  d_\Sigma=\mathsf H^{-1}(\Sigma/\varepsilon),
  \qquad
  \delta a_\sigma=\frac{(d_\Sigma)_\sigma}{2a_\sigma},
  \qquad
  \mathcal L\Sigma=\sqrt{\varepsilon}\sum_{\sigma=\pm}\delta a_\sigma b_\sigma.
  \tag{7.31}\label{eq:7.31}
\end{equation}
The divisions are componentwise and are defined first on
\NSi{NS-F-P084-012}{X_a<X<X_b}, where \NSi{NS-F-P084-013}{a_\sigma>0}. The
estimates below justify smooth zero extension at the shell edges.

\begin{proposition}\label{prop:7.6}
For every \NSi{NS-F-P084-014}{\alpha\in\mathbb R} and every real two-vector
\NSi{NS-F-P084-015}{\Sigma(R,Z,T)\in\mathsf M_\alpha} independent of the angular and
auxiliary variables, \eqref{eq:7.31} defines a linear operation
\NSi{NS-F-P084-016}{\mathcal L\Sigma}. Its harmonic coefficients on the slow
neighborhood satisfy the local derivative and envelope bounds of
\NSi{NS-F-P084-017}{\mathsf W_{\alpha-1/2}}, with smooth zero extension at the shell
edges. With the slow cutoffs shown explicitly, its bounds and covariance are
\NSFormula{NS-F-P084-018}{display}{7.32}
\begin{equation}
  \eta_\beta\mathcal L\Sigma\in\mathsf W_{\alpha-1/2},
  \qquad
  \mathcal B(W_0,\mathcal L\Sigma)=\Sigma,
  \qquad
  \mathcal C(\eta_\beta\mathcal L\Sigma)\in\mathsf M_{2\alpha-1}.
  \tag{7.32}\label{eq:7.32}
\end{equation}
The denominator \NSi{NS-F-P084-019}{a_\sigma} remains this fixed coefficient
at every correction stage.
\end{proposition}

No sign or relative-size restriction is imposed on
\NSi{NS-F-P084-020}{\Sigma} in Proposition~\ref{prop:7.6}.

\begin{proof}
Apply Leibniz' rule to
\NSi{NS-F-P084-021}{d_\Sigma=\mathsf H^{-1}(\Sigma/\varepsilon)}, using
\eqref{eq:7.29} and the differentiated inverse bounds proved in
Proposition~\ref{prop:7.5}. The weighted derivative bounds on
\NSi{NS-F-P084-022}{\Sigma} give
\NSi{NS-F-P084-023}{|D^I(d_\Sigma)_\sigma|
\leq C_I\varepsilon^{\alpha-1}S_*^{b'_I}\zeta\delta^{-a'_I}}. Every derivative
of \NSi{NS-F-P084-024}{y_\sigma^{-1/2}} is a sum of terms of the form
\NSi{NS-F-P084-025}{Cy_\sigma^{-k-1/2}
\prod_{\nu=1}^kD^{I_\nu}y_\sigma}. The \NSi{NS-F-P084-026}{k} numerator
weights cancel all but \NSi{NS-F-P084-027}{\zeta^{-1/2}} in the denominator,
by \eqref{eq:7.25}. Leibniz therefore gives
\NSFormula{NS-F-P084-028}{display}{7.33}
\begin{equation}
  |D^I\delta a_\sigma|
  \leq C_I\varepsilon^{\alpha-1}S_*^{b'_I}\sqrt{\zeta}\,\delta^{-a'_I}.
  \tag{7.33}\label{eq:7.33}
\end{equation}
The remaining square-root edge weight proves smooth zero extension;
Lemma~\ref{lem:7.4} gives the local wave estimates. Multiplication by
\NSi{NS-F-P084-029}{\eta_\beta} gives the supported wave class. Write
\NSi{NS-F-P084-030}{a=(a_+,a_-)^T} and
\NSi{NS-F-P084-031}{\delta a=(\delta a_+,\delta a_-)^T}, and let
\NSi{NS-F-P084-032}{\odot} denote componentwise multiplication. Since these
scalar coefficients are independent of the averaging variables, disjoint
rectangles give the exact bilinear identity
\NSFormula{NS-F-P084-033}{display}{7.34}
\begin{equation}
  \mathcal B(W_0,\mathcal L\Sigma)
  =\varepsilon\mathsf H(2a\odot\delta a)=\varepsilon\mathsf Hd_\Sigma=\Sigma.
  \tag{7.34}\label{eq:7.34}
\end{equation}
The factor \NSi{NS-F-P084-034}{1/2} for the real cosine has already entered
\eqref{eq:7.27}; the factor \NSi{NS-F-P084-035}{2} in this cross identity
explains the denominator in \eqref{eq:7.31}. In
\NSi{NS-F-P084-036}{\mathcal C(\eta_\beta\mathcal L\Sigma)}, both wave factors
have exponent \NSi{NS-F-P084-037}{\alpha-1/2}, so the class product rule gives
the last assertion in \eqref{eq:7.32}.
\end{proof}
